%% toc: P2 14 14.1,14.2,14.3,14.4
\section{Proofs for finite inputs and the arbitrary-language theorem}
\label{sec:annex-p02-w02-p02-016-finite-input-scope}

Throughout, \(P\) is a set, \(\operatorname{Form}(P)\), valuations,
\(\entails_P\), \(\equiv_P\) and \(\operatorname{Var}\) are as in
Convention~\ref{def:p02-w02-p02-016-semantics}, and \(\mathsf r^{\star}\) is the
translation of Definition~\ref{def:p02-w02-p02-016-renaming}.  All arguments
take place in \(\mathsf{ZF}\).  Where a finite set is used, its finiteness
gives a bijection with some \(\{0,\ldots,k-1\}\), and only that one bijection
is fixed; no family of choices is made.

\begin{annexitem}{annex:p02-w02-p02-016-atoms}{thm:p02-w02-p02-016-finite-atoms}{Finitely many atoms}
\begin{proof}
(1) By structural induction on \(\chi\).  For an atom \(a\),
\(\operatorname{Var}(a)=\{a\}\) is finite and \(a\) has no immediate
subformula.  For \(\bottom\), \(\operatorname{Var}(\bottom)=\varnothing\).
For \(\neg\chi\), \(\operatorname{Var}(\neg\chi)=\operatorname{Var}(\chi)\),
which is finite by the induction hypothesis, and the only immediate
subformula is \(\chi\).  For \(\chi\circ\theta\),
\(\operatorname{Var}(\chi\circ\theta)=
\operatorname{Var}(\chi)\cup\operatorname{Var}(\theta)\) contains the
variable sets of both immediate subformulas.  It is finite: by the induction
hypothesis there are bijections \(f:\operatorname{Var}(\chi)\to\{0,\ldots,m-1\}\)
and \(g:\operatorname{Var}(\theta)\to\{0,\ldots,n-1\}\); the map sending \(x\)
to \(f(x)\) when \(x\in\operatorname{Var}(\chi)\) and to \(m+g(x)\) otherwise
is an injection of the union into \(\{0,\ldots,m+n-1\}\), and every subset of
a finite initial segment of \(\N\) is finite (by induction on the length of
the segment).  Only the two bijections for the two immediate subformulas are
used at this step.

(2) By induction on the number \(n\) of members of \(\Gamma\).  If \(n=0\),
\(\operatorname{Var}(\Gamma)=\varnothing\).  If \(\Gamma\) has \(n+1\)
members, fix a bijection \(\{0,\ldots,n\}\to\Gamma\) and let \(\gamma\) be the
image of \(n\) and \(\Gamma'\) the image of the rest; then
\(\operatorname{Var}(\Gamma)=\operatorname{Var}(\Gamma')\cup
\operatorname{Var}(\gamma)\), a union of two finite sets by the induction
hypothesis and (1), which is finite by the argument in (1).

(3) Each atom \(a\in P\), as an atomic formula, has
\(\operatorname{Var}(a)=\{a\}\), so
\(\operatorname{Var}(P)=\bigcup_{a\in P}\{a\}=P\).  This set is infinite
whenever \(P\) is.
\end{proof}
\end{annexitem}

\begin{annexitem}{annex:p02-w02-p02-016-transfer}{thm:p02-w02-p02-016-transfer}{Finite-input transfer}
\begin{proof}
(1) The set \(\Gamma\cup\{\varphi,\psi\}\) is finite, so \(A\) is finite by
Lemma~\ref{thm:p02-w02-p02-016-finite-atoms}(2).  Fix a bijection
\(\beta:A\to\{0,\ldots,k-1\}\), where \(k\) is the number of elements of
\(A\), and put \(\mathsf r(a)=p_{\beta(a)}\).  Since \(n\mapsto p_n\) is injective,
\(\mathsf r\) is a bijection onto \(\{p_0,\ldots,p_{k-1}\}\), so it is initial
for \(A\).

For (2) fix any injection \(\mathsf r:A\to\mathsf{Prop}\).  Two transfer claims
connect the valuations on the two sides.

\emph{Claim A (pulling back).}  For a valuation
\(w:\mathsf{Prop}\to\{\mathbf T,\mathbf F\}\), define
\(w^{\mathsf r}:P\to\{0,1\}\) by \(w^{\mathsf r}(a)=1\) iff
\(\widehat w(\mathsf r^{\star}(a))=\mathbf T\).  Then for every
\(\chi\in\operatorname{Form}(P)\),
\(\widehat{w^{\mathsf r}}(\chi)=1\) iff \(\widehat w(\mathsf r^{\star}(\chi))=\mathbf T\).
By structural induction on \(\chi\).  For an atom this is the definition of
\(w^{\mathsf r}\).  For \(\bottom\), \(\mathsf r^{\star}(\bottom)=\bottom\) and both
values are false.  For \(\neg\chi\):
\(\widehat{w^{\mathsf r}}(\neg\chi)=1\) iff \(\widehat{w^{\mathsf r}}(\chi)=0\) iff (by
the hypothesis) \(\widehat w(\mathsf r^{\star}(\chi))=\mathbf F\) iff
\(\widehat w(\neg\mathsf r^{\star}(\chi))=\mathbf T\), and
\(\neg\mathsf r^{\star}(\chi)=\mathsf r^{\star}(\neg\chi)\).  For \(\chi\circ\theta\),
the truth clause for \(\circ\) computes both sides from the values of the two
immediate subformulas on each side, which correspond by the hypothesis, and
\(\mathsf r^{\star}(\chi\circ\theta)=(\mathsf r^{\star}(\chi)\circ\mathsf r^{\star}(\theta))\).

\emph{Claim B (pushing forward).}  For a valuation \(v:P\to\{0,1\}\), define
\(v_{\mathsf r}:\mathsf{Prop}\to\{\mathbf T,\mathbf F\}\) by
\(v_{\mathsf r}(p_n)=\mathbf T\) iff \(p_n=\mathsf r(a)\) for some \(a\in A\) with
\(v(a)=1\).  Because \(\mathsf r\) is injective, at most one \(a\) has
\(\mathsf r(a)=p_n\), so \(v_{\mathsf r}(\mathsf r(a))=\mathbf T\) iff \(v(a)=1\) for
\(a\in A\).  Then for every \(\chi\) with
\(\operatorname{Var}(\chi)\subseteq A\),
\(\widehat{v_{\mathsf r}}(\mathsf r^{\star}(\chi))=\mathbf T\) iff
\(\widehat v(\chi)=1\).  By structural induction over the formulas with
variables in \(A\); this class is closed under immediate subformulas by
Lemma~\ref{thm:p02-w02-p02-016-finite-atoms}(1).  For an atom \(a\in A\),
\(\mathsf r^{\star}(a)=\mathsf r(a)\) and the equivalence was just noted.  The
\(\bottom\) and connective cases are those of Claim~A.

\emph{Consequence.}  Suppose \(\Gamma\entails_P\varphi\) and let \(w\)
satisfy \(\mathsf r^{\star}[\Gamma]\).  By Claim~A, \(w^{\mathsf r}\) satisfies every
member of \(\Gamma\), hence satisfies \(\varphi\), and by Claim~A again
\(\widehat w(\mathsf r^{\star}(\varphi))=\mathbf T\).  Conversely suppose
\(\mathsf r^{\star}[\Gamma]\entails\mathsf r^{\star}(\varphi)\) and let \(v\) satisfy
\(\Gamma\).  Every member of \(\Gamma\) has its variables in \(A\), so by
Claim~B \(v_{\mathsf r}\) satisfies \(\mathsf r^{\star}[\Gamma]\), hence satisfies
\(\mathsf r^{\star}(\varphi)\); since \(\operatorname{Var}(\varphi)\subseteq A\),
Claim~B gives \(\widehat v(\varphi)=1\).

\emph{Satisfiability.}  If \(v\) satisfies \(\varphi\), then \(v_{\mathsf r}\)
satisfies \(\mathsf r^{\star}(\varphi)\) by Claim~B.  If \(w\) satisfies
\(\mathsf r^{\star}(\varphi)\), then \(w^{\mathsf r}\) satisfies \(\varphi\) by
Claim~A.

\emph{Validity.}  Validity over \(P\) is \(\varnothing\entails_P\varphi\), and
validity in \(\mathsf{Fm}_{\mathsf{CPL}}\) is
\(\varnothing\entails\mathsf r^{\star}(\varphi)\).  The argument for consequence
used only \(\operatorname{Var}(\Gamma\cup\{\varphi\})\subseteq A\), which holds
for \(\varnothing\) in place of \(\Gamma\).

\emph{Equivalence.}  Suppose \(\varphi\equiv_P\psi\).  For every \(w\),
Claim~A gives \(\widehat w(\mathsf r^{\star}(\varphi))=\mathbf T\) iff
\(\widehat{w^{\mathsf r}}(\varphi)=1\) iff \(\widehat{w^{\mathsf r}}(\psi)=1\) iff
\(\widehat w(\mathsf r^{\star}(\psi))=\mathbf T\).  Conversely suppose
\(\mathsf r^{\star}(\varphi)\equiv\mathsf r^{\star}(\psi)\).  For every \(v\), since
\(\operatorname{Var}(\varphi)\cup\operatorname{Var}(\psi)\subseteq A\),
Claim~B gives \(\widehat v(\varphi)=1\) iff
\(\widehat{v_{\mathsf r}}(\mathsf r^{\star}(\varphi))=\mathbf T\) iff
\(\widehat{v_{\mathsf r}}(\mathsf r^{\star}(\psi))=\mathbf T\) iff
\(\widehat v(\psi)=1\).

(3) In each equivalence of (2) the left side does not mention \(\mathsf r\), so
the right side has the same truth value for every injection
\(\mathsf r:A\to\mathsf{Prop}\).

Every map above is defined explicitly from the one fixed \(\mathsf r\) (and, in
(1), the one fixed \(\beta\)); every induction is structural or on a natural
number.  No choice principle, no Boolean prime ideal, and no completeness or
compactness theorem enters.
\end{proof}
\end{annexitem}

\begin{annexitem}{annex:p02-w02-p02-016-decision}{thm:p02-w02-p02-016-decision}{Finite inputs over an arbitrary set are settled in the countable coding}
\begin{proof}
By Definition~\ref{def:p02-w02-p02-016-renaming}, \(\mathsf r^{\star}\) commutes
with every constructor, so
\(\mathsf r^{\star}((\bigwedge\Gamma)\to\varphi)=
(\bigwedge\mathsf r^{\star}[\Gamma]\to\mathsf r^{\star}(\varphi))\), where
\(\bigwedge\mathsf r^{\star}[\Gamma]\) is the right-associated conjunction of
\(\mathsf r^{\star}(\gamma_1),\ldots,\mathsf r^{\star}(\gamma_m)\), and
\(\mathsf r^{\star}(\topv)=\mathsf r^{\star}(\neg\bottom)=\neg\bottom\).

\emph{Claim C.}  For every \(w:\mathsf{Prop}\to\{\mathbf T,\mathbf F\}\),
\(\widehat w(\bigwedge\mathsf r^{\star}[\Gamma])=\mathbf T\) iff \(w\) makes every
\(\mathsf r^{\star}(\gamma_i)\) true.  By induction on \(m\): for \(m=0\),
\(\neg\bottom\) is true under every valuation and the empty condition holds;
for \(m=1\) the conjunction is \(\mathsf r^{\star}(\gamma_1)\); for \(m\ge2\) it
is \((\mathsf r^{\star}(\gamma_1)\land\bigwedge(\mathsf r^{\star}(\gamma_2),\ldots,
\mathsf r^{\star}(\gamma_m)))\), and the truth clause for \(\land\) with the
induction hypothesis gives the claim.

(1) By Theorem~\ref{thm:p02-w02-p02-016-transfer}(2),
\(\Gamma\entails_P\varphi\) iff
\(\mathsf r^{\star}[\Gamma]\entails\mathsf r^{\star}(\varphi)\).  The latter holds iff
\(\theta=(\bigwedge\mathsf r^{\star}[\Gamma]\to\mathsf r^{\star}(\varphi))\) is valid:
under a valuation \(w\) that satisfies \(\mathsf r^{\star}[\Gamma]\), Claim~C makes
the antecedent true, so \(\theta\) is true iff \(\mathsf r^{\star}(\varphi)\) is;
under any other \(w\), Claim~C makes the antecedent false and \(\theta\) true.
By Definition~\ref{def:p02-sat-taut}, \(\mathsf{TAUT}\) is the set of
encodings of valid formulas, and the encoding has an inverse, so the encoding
of \(\theta\) lies in \(\mathsf{TAUT}\) iff \(\theta\) is valid.

(2) The satisfiability and validity clauses are
Theorem~\ref{thm:p02-w02-p02-016-transfer}(2) combined with the definitions of
\(\mathsf{SAT}\) and \(\mathsf{TAUT}\) in Definition~\ref{def:p02-sat-taut},
again using that the encoding is invertible.  For equivalence,
\(\mathsf r^{\star}((\varphi\to\psi)\land(\psi\to\varphi))=
((\mathsf r^{\star}(\varphi)\to\mathsf r^{\star}(\psi))\land
(\mathsf r^{\star}(\psi)\to\mathsf r^{\star}(\varphi)))\), and by the truth clauses for
\(\to\) and \(\land\) this formula is true under \(w\) iff
\(\mathsf r^{\star}(\varphi)\) and \(\mathsf r^{\star}(\psi)\) receive the same value
under \(w\).  So it is valid iff
\(\mathsf r^{\star}(\varphi)\equiv\mathsf r^{\star}(\psi)\), which by
Theorem~\ref{thm:p02-w02-p02-016-transfer}(2) holds iff \(\varphi\equiv_P\psi\).

(3) First, \(\operatorname{Var}(\mathsf r^{\star}(\chi))=\mathsf r[\operatorname{Var}(\chi)]\)
for every \(\chi\) with \(\operatorname{Var}(\chi)\subseteq A\), by
structural induction: an atom \(a\in A\) gives \(\{\mathsf r(a)\}\), \(\bottom\)
gives \(\varnothing\) on both sides, and the connective cases are unions on
both sides.  Each formula named in (1) and (2) is \(\mathsf r^{\star}\) of a
formula whose variables lie in \(A\) (for \(\topv\) there are none), so its
variables lie in \(\mathsf r[A]=X\), where \(X=\{p_0,\ldots,p_{k-1}\}\) because
\(\mathsf r\) is initial.  Let \(\chi'\) be such a formula.  Every valuation
\(w\) restricts to a function \(u=w\restriction X\), and every function
\(u:X\to\{\mathbf T,\mathbf F\}\) is the restriction of the valuation \(w_u\)
that agrees with \(u\) on \(X\) and gives \(\mathbf F\) to every other
variable.  Since \(w\) and \(w_{w\restriction X}\) agree on every variable
occurring in \(\chi'\), Lemma~\ref{lem:prop-coincidence} gives
\(\widehat w(\chi')=\widehat{w_{w\restriction X}}(\chi')\).  Hence \(\chi'\) is
satisfiable iff \(\widehat{w_u}(\chi')=\mathbf T\) for some \(u\), and valid iff
this holds for every \(u\).  There are exactly \(2^k\) functions
\(X\to\{\mathbf T,\mathbf F\}\), since \(X\) has \(k\) elements.
\end{proof}
\end{annexitem}

\begin{annexitem}{annex:p02-w02-p02-016-no-coding}{thm:p02-w02-p02-016-no-coding}{No encoding or enumeration of an uncountable formula set}
\begin{proof}
Let \(j:P\to\operatorname{Form}(P)\) send each atom to the atomic formula it
forms; \(j\) is injective because distinct atoms are distinct formulas
(Convention~\ref{def:p02-fbc-language} keeps atoms apart from every other
token).

(1) Define \(\pi:\{0,1\}^*\to\N\) by
\[
 \pi(x_1\cdots x_n)=2^{n}-1+\sum_{i=1}^{n}x_i\,2^{\,n-i}.
\]
For fixed \(n\), the sum is the number with binary digits \(x_1\cdots x_n\),
and distinct digit strings of length \(n\) give distinct numbers between
\(0\) and \(2^n-1\); so \(\pi\) is injective on \(\{0,1\}^n\) with values
between \(2^n-1\) and \(2^{n+1}-2\).  These ranges are disjoint for distinct
\(n\), so \(\pi\) is injective.  If \(e:\operatorname{Form}(P)\to\{0,1\}^*\)
were injective, then \(\pi\circ e\circ j:P\to\N\) would be injective,
contradicting the hypothesis on \(P\).  An encoding of a set of formulas by
bit strings that has an inverse, as in Definition~\ref{def:p02-sat-taut}, is an
injection into \(\{0,1\}^*\); so no such encoding of \(\operatorname{Form}(P)\)
exists.  In particular \(\operatorname{Form}(P)\) is not the set of formulas
encoded by Definition~\ref{def:p02-sat-taut}, which is
\(\mathsf{Fm}_{\mathsf{CPL}}\) and is injectively encoded.

(2) Suppose \(s:\N\to\operatorname{Form}(P)\) is surjective.  For each
\(\chi\), the set \(\{n:s(n)=\chi\}\) is nonempty, and its least member
\(g(\chi)\) exists because \(\N\) is well-ordered; no choice is involved.
Since \(s(g(\chi))=\chi\), the map \(g\) is injective, and \(g\circ j:P\to\N\)
is injective, again contradicting the hypothesis.
\end{proof}
\end{annexitem}
